\section{Notation and Minimal Assumptions}


本稿で用いる記号は、局所的な作業記法であり、既存理論の標準的意味をそのまま引き継ぐものではない。

\begin{itemize}
\item $t$: 外部から測定可能な物理時間。
\item $\tau$: 内部時間遅延、バッファリング、保持。主観的時間そのものではない。
\item $v_f$: 高速な予測的・反応的ダイナミクス。
\item $v_s$: 低速な沈殿的・安定化的ダイナミクス。
\item $\Delta v = v_f - v_s$: 速度差。
\item $\mu$: 速度スケール間の結合粘性。
\item $\mathcal{T}$: 速度差を変換する有効なトルク様量。
\item $\Phi$: 情報ポテンシャル。
\item $I$: 情報密度。
\item $J$: 情報流。
\end{itemize}

ここで重要なのは、$\tau$ をトルクとして使わないことである。$\tau$ は内部遅延であり、トルク様量は一貫して $\mathcal{T}$ と表す。両者はいずれも構造式の中で $\Delta v/\mu$ を含むが、記述している側面は異なる。$\tau$ は参照の時間的厚みであり、$\mathcal{T}$ は速度差の変換負荷である。

これらの構造関係は第7節で整理する。
